\subsection{The remaining finite-field cusp fibre}

Let \(a\in\{5,7\}\), let \(X_a\) be the proper Fermat curve, and let
\(D_a\) be its coordinate boundary. Suppose a full geometrically
pro-\(3\) point-or-cusp Selmer section has compact image equal to that
of a rational coordinate cusp \(P\). For \(p\notin\{3,a\}\), put
\[
 E_{a,P}(p)=
 \left\{Q\in D_a(\mathbf F_p):
 [s_Q]=[s_P]\text{ in }
 \operatorname{Sec}\bigl(\pi_1^{\mathrm{geom},(3)}
             (X_a/\mathbf F_p)\bigr)\right\}.
\]
Section classes here are taken modulo the geometric pro-\(3\) group;
the arithmetic quotient remains \(G_{\mathbf F_p}\).
Smooth proper specialization shows that every centre of nontrivial
local affine inertia belongs to \(E_{a,P}(p)\). Indeed, a local cusp,
or a local point reducing to a cusp \(Q\), has compact specialization
\(s_Q\); a point with interior reduction is unramified at \(p\).
For a section defined only in a geometric quotient \(R\), this
conclusion uses its compact quotient \(R^{\mathrm{cp}}\), and the
corresponding fibre can be larger. No full global lift is implicit.

General point injectivity for a single geometric pro-\(\ell\) quotient
over a finite field is false. Hoshi's Remark~10(i), which records an
example of Tamagawa, gives distinct points with the same full
geometrically pro-\(\ell\) section class
\cite[pp.~742--743]{FCHoshi2009}. After a finite constant extension,
the arithmetic group has the form
\[
 \Pi_X\simeq A\times \prod_{q\ne\ell}\mathbf Z_q,\qquad
 1\longrightarrow\Delta_X\longrightarrow A
       \longrightarrow\mathbf Z_\ell\longrightarrow1.
\]
For a further constant extension \(L/K\) of degree prime to \(\ell\),
new \(L\)-points conjugate under \(\operatorname{Gal}(L/K)\) have the
same section class: conjugation by a constant-field lift differs from
conjugation inside \(\Pi_{X_L}\) by a central factor. Since \(G_L\) is
abelian, conjugacy by \(\Pi_{X_L}\) is equivalent to geometric conjugacy
for sections. This example does not identify the points as Fermat
coordinate cusps, does not involve a prescribed rational cusp \(P\),
and does not compute \(E_{a,P}(p)\) over a prime field.

The almost-injectivity theorem of Saidi--Tamagawa requires a cofinite
set \(\Sigma\) of retained geometric primes
\cite[Proposition~1.8(vi)--(vii), Corollary~1.11(iii)]
 {FCSaidiTamagawa2009}.
It cannot be applied with \(\Sigma=\{3\}\). Its affine inertia
statements also do not restore point separation after boundary inertia
is killed in the proper quotient.

Even a proof that \(E_{a,P}(p)=\{P\}\) on a density-one set of primes
would not by itself prove finite ramification support for the
transverse boundary class. It would only make that class unramified on
the density-one set; the complement may be infinite. It would not make
the local class zero, since an interior reduction need not belong to
the \(P\)-cusp packet. The unrestricted cusp parameter itself is
allowed to have infinitely many nonzero valuations, so the needed
support assertion concerns the transverse component.

\begin{lemma}[Density-one Kummer localization]\label{lem:density-one-kummer-localization}
Let \(F\) be a number field. Localization on any density-one set of its
finite places, in Dirichlet density, is injective on \(H^1(F,\mathbf Z_3(1))\). The analogous
statement over \(\mathbf Q\) holds for finite direct sums of induced
Tate modules
\(\operatorname{Ind}_{G_F}^{G_{\mathbf Q}}\mathbf Z_3(1)\).
\end{lemma}
\begin{proof}
For \(C=\mu_{3^m}\) and \(E=F(C)\), the cyclic group
\(G=\operatorname{Gal}(E/F)\subset(\mathbf Z/3^m)^\times\) satisfies
\(H^1(G,C)=0\). If a generator \(u\) is nontrivial modulo \(3\), then
\(u-1\) is invertible. Otherwise, writing
\(v_3(u-1)=s<m\), the order is \(3^{m-s}\), the norm has valuation
\(m-s\), and its kernel is \(3^sC=(u-1)C\). The trivial group is
immediate. Restriction to \(E\) is therefore injective.

A nonzero finite-level class consequently defines a nontrivial
Kummer extension over \(E\). A nonidentity element of its Kummer
subgroup, viewed in the Galois splitting field over \(F\), is detected
by Frobenius at a positive-density set of primes splitting in \(E\).
At such a prime the coefficient action is trivial and the local
class is nonzero. This contradicts density-one local vanishing.
The inverse systems of finite zeroth cohomology groups satisfy the Mittag--Leffler condition, so passage through all \(m\) proves the first assertion. Shapiro's lemma
proves the second: the sum of $(N\mathfrak p)^{-s}$ over primes of residue degree at least
two remains bounded as $s\to1^+$, while the excluded degree-one prime
sum is at most $[F:\mathbf Q]$ times the excluded rational prime sum.
Thus lifting a rational Dirichlet-density-one set gives a
Dirichlet-density-one set of places of $F$.
\end{proof}

Thus density-one \emph{local vanishing} of the transverse boundary
class would suffice, whereas density-one separation of coordinate
cusps supplies only an inertia restriction. No singleton computation
of the above Fermat fibres, finite-support theorem for the transverse
class, or affine pro-\(3\) classification is asserted here.
